\section{Khi có lỗi}
\label{sec:loi}

Gặp lỗi trong quá trình soạn thảo là điều hoàn toàn bình thường, và SumiTeX
cung cấp sẵn các công cụ giúp bạn phát hiện cũng như khắc phục chúng nhanh
chóng.

\subsection{Lỗi ngay khi gõ}
Nhiều lỗi cơ bản sẽ được phát hiện ngay trong lúc bạn gõ, trước cả khi biên
dịch. Khi đó, một biểu tượng ⚠ sẽ xuất hiện ở đầu dòng kèm theo lời giải thích
ngắn gọn ở cuối dòng.

\anh[0.85\linewidth]{loi-goi}{Lệnh \lenh{mathbb} yêu cầu phải khai báo gói
\texttt{amssymb}}

Bạn chỉ cần đặt con trỏ tại dòng báo lỗi rồi nhấn \phim{Ctrl}+\phim{.}, SumiTeX
sẽ tự động thêm gói còn thiếu vào phần đầu tài liệu.

\subsection{Bảng lỗi}
Sau khi biên dịch, bạn có thể mở bảng lỗi bằng cách nhấp vào đèn báo lỗi hoặc
nhấn \phim{Ctrl}+\phim{Shift}+\phim{M}.

\anh[0.75\linewidth]{van-de}{Bảng lỗi: mỗi dòng tương ứng với một lỗi, nhấp vào
để nhảy thẳng tới vị trí sai}

\begin{meo}
\textbf{Hãy luôn sửa lỗi đầu tiên trước.} Một lỗi nhỏ ở trên thường kéo theo
hàng loạt lỗi dây chuyền ở phía dưới; do đó, khi bạn sửa xong lỗi đầu tiên và
biên dịch lại, phần lớn các lỗi còn lại sẽ tự biến mất. Ngoài ra, phím \phim{F8}
giúp bạn di chuyển nhanh tới lỗi kế tiếp.
\end{meo}

\subsection{Các lỗi thường gặp}
\begin{description}
  \item[Unicode character …] Tài liệu chưa được cấu hình tiếng Việt. Bạn hãy
    thêm đủ ba dòng \lenh{usepackage} như hướng dẫn ở chương~\ref{sec:van-ban}.
  \item[Undefined control sequence] Bạn đã gõ sai tên lệnh hoặc quên khai báo
    gói thư viện tương ứng. Hãy xem gợi ý sửa lỗi ở cuối dòng.
  \item[Missing \$ inserted] Bạn đang dùng ký hiệu toán học (như \verb|^| hoặc
    \verb|_|) ở ngoài môi trường toán. Hãy đặt biểu thức đó vào giữa hai dấu
    \texttt{\$}.
  \item[Missing \} inserted] Câu lệnh bị thiếu dấu đóng ngoặc nhọn \}.
\end{description}

\begin{thuchanh}[Bài 6]
Bài 6 được cố ý tạo sẵn một số lỗi nên đang tạm tắt. Hãy mở tệp
\texttt{main.tex}, xóa dấu \texttt{\%} ở đầu dòng có lệnh
\texttt{\textbackslash input\{bai-06\}}, sau đó biên dịch và sửa dần từng lỗi
cho đến khi bảng lỗi hoàn toàn trống.
\end{thuchanh}